% ==========================================================================
%  related_work.tex — Physically-Guided Collision-Free Indoor Scene Generation
% ==========================================================================
\section{Related Work}
\label{sec:related}

% ------------------------------------------------------------------
\subsection{3D Indoor Scene Synthesis}
\label{sec:related:scene_synth}
% ------------------------------------------------------------------

Indoor scene synthesis aims to arrange furniture and objects within a room in configurations that are both aesthetically pleasing and functionally plausible.
Early approaches relied on handcrafted rules, design guidelines, and optimization over spatial grammars~\cite{placeholder_grammars, placeholder_yeh2012, placeholder_fisher2012}.
Data-driven methods subsequently learned arrangement priors from real-world datasets, formulating layout generation as autoregressive sequence prediction~\cite{placeholder_autoreg_ritchie, placeholder_sceneformer, placeholder_atiss} or variational inference over object attributes~\cite{placeholder_fastsynth, placeholder_planit}.
With the advent of diffusion models~\cite{placeholder_ddpm_ho, placeholder_song_score}, recent work has shown that denoising-based generation can produce diverse, high-quality layouts by iteratively refining object positions, sizes, and orientations from noise~\cite{placeholder_diffuscene, placeholder_layoutdiffusion}.
While these methods achieve strong distributional fidelity, they generally optimize reconstruction losses that do not explicitly penalize inter-object collisions, resulting in scenes that are statistically plausible but often physically invalid.

% ------------------------------------------------------------------
\subsection{Scene Graph-Conditioned Generation}
\label{sec:related:sg_gen}
% ------------------------------------------------------------------

Scene graphs provide a structured, human-interpretable interface for controllable generation by encoding object categories as nodes and spatial relationships as labeled edges.
In 2D, scene graph-to-image synthesis has been widely studied through adversarial~\cite{placeholder_sg2im, placeholder_canonical_sg2im} and diffusion-based~\cite{placeholder_sgdiff_2d} approaches.
Extending this paradigm to 3D, Graph-to-3D~\cite{placeholder_graph2threed} proposed a variational framework that generates object shapes and layouts from scene graphs.
CommonScenes~\cite{placeholder_commonscenes} improved upon this by collapsing the graph into subject--predicate--object triplets and applying graph convolution to aggregate per-node features, which then condition a latent diffusion model for shape generation.
However, CommonScenes processes each object's denoising independently, leading to limited global coherence across the generated scene.

EchoScene~\cite{echoscene} addressed this limitation through a dual-branch diffusion framework (layout and shape) equipped with an \emph{Information Exchange Unit}---a triplet-GCN that, at every denoising timestep, aggregates the current noisy states of neighboring nodes via the original graph edges.
This ``echo'' mechanism ensures each object's denoiser is aware of its neighbors' evolving states, improving style consistency and manipulation quality.
Our work builds upon EchoScene's layout branch and complements its relational awareness with explicit physical constraints, addressing the collision problem that the echo mechanism alone cannot guarantee to solve.

% ------------------------------------------------------------------
\subsection{Physical Plausibility in Scene Generation}
\label{sec:related:physical}
% ------------------------------------------------------------------

Ensuring physical plausibility---particularly collision avoidance---in generated 3D scenes has received growing attention.
One line of work incorporates physics constraints during \emph{training}.
ScenePriors~\cite{placeholder_scenepriors} learns a layout prior via normalizing flows and applies collision penalties as regularization during the variational learning process.
DiffuScene~\cite{placeholder_diffuscene} integrates scene-level constraints into the diffusion training objective.
However, training-time approaches require access to the full training pipeline and risk altering the learned distribution in unpredictable ways.

A second line of work applies constraints at \emph{inference time} or during \emph{post-processing}.
PhyScene~\cite{placeholder_physcene} pioneered inference-time physical guidance for layout diffusion, introducing differentiable losses for collision avoidance, room boundary compliance, and walkability (path reachability) that steer the reverse sampling process via gradient updates.
LEGO-Net~\cite{placeholder_legonet} learns to predict physically valid object placements through an iterative refinement network.
ControlRoom3D~\cite{placeholder_controlroom3d} uses a mesh-level collision detection and response to refine generated rooms.

Our work draws inspiration from PhyScene's inference-time guidance paradigm but differs in several important ways.
First, we integrate the physical guidance into a \emph{scene graph-conditioned} diffusion model (EchoScene), whereas PhyScene operates on unconditional or text-conditioned layout diffusion---the scene graph structure introduces additional challenges such as auxiliary non-object nodes (floor, scene boundaries) that must be handled via objectness masking and ground-truth freezing.
Second, we introduce a complementary \emph{post-processing stage} based on the Separating Axis Theorem that provides geometrically exact, rotation-aware collision resolution.
This coarse-to-fine decomposition is novel: to our knowledge, no prior work combines inference-time physical guidance with a deterministic geometric post-processor for scene graph-conditioned generation.

% ------------------------------------------------------------------
\subsection{Collision Detection and Resolution}
\label{sec:related:collision}
% ------------------------------------------------------------------

Collision detection is a classic problem in computational geometry and real-time simulation~\cite{placeholder_ericson_rtcd}.
For convex polygons, the Separating Axis Theorem (SAT)~\cite{placeholder_sat_gottschalk} provides an exact and efficient test: two convex shapes are separated if and only if there exists an axis along which their projections do not overlap.
For pairs of 2D rectangles (or 3D cuboids projected onto a plane), SAT requires testing only four candidate axes (the edge normals of each rectangle), making it $O(1)$ per pair.

In the context of scene generation, prior collision resolution methods typically operate on axis-aligned bounding boxes (AABBs)~\cite{placeholder_aabb_methods}, which ignore object rotation and can both miss true collisions (false negatives for rotated objects) and report spurious overlaps (false positives for diagonal arrangements).
GIoU-based penalties~\cite{placeholder_giou} offer differentiable alternatives but remain axis-aligned.
More recently, differentiable oriented IoU computations~\cite{placeholder_oriented_iou_zhou} have been proposed, but these are designed for loss computation during training rather than exact post-hoc resolution.

Our post-processing module uses SAT on oriented bounding boxes in the floor plane, providing exact collision detection that accounts for yaw rotation.
Unlike standard minimum-translation-vector (MTV) resolution, which can push objects in unintuitive directions (e.g., behind an adjacent piece of furniture), we employ a \emph{center-biased axis selection} strategy that preferentially pushes objects apart along the direction connecting their centers, producing more natural separations.

% ------------------------------------------------------------------
\subsection{Classifier Guidance for Diffusion Models}
\label{sec:related:classifier_guidance}
% ------------------------------------------------------------------

Our physics-guided diffusion mechanism is conceptually related to classifier guidance~\cite{placeholder_classifier_guidance_dhariwal} and its generalizations.
In the original formulation, a separately trained classifier provides gradients that steer the reverse process toward a target class.
Subsequent work has extended this idea to arbitrary differentiable objectives: DPS~\cite{placeholder_dps} uses measurement-consistency losses for inverse problems, and Universal Guidance~\cite{placeholder_universal_guidance} demonstrates that any differentiable loss function can guide diffusion sampling.
In the 3D domain, DreamFusion~\cite{placeholder_dreamfusion} uses Score Distillation Sampling to guide 3D representations via a 2D diffusion prior.

We apply this principle to scene layout generation: our physical guidance losses (collision, room boundary, reachability) serve as the ``classifier,'' and their gradients are injected into the reverse process via variance-preconditioned mean shifting.
The key design consideration specific to our setting is \emph{scheduling}---we delay guidance activation until the diffusion process has established a coarse layout, avoiding destabilization of early high-noise steps where the predicted clean sample $\hat{\mathbf{d}}_0^{(t)}$ is unreliable.
